\clearpage
\VerdierPhysicalPageAnchor{0184}%
\begingroup\small\setlength{\parskip}{0pt}\setlength{\emergencystretch}{1.5em}
\setlength{\abovedisplayskip}{4pt plus 2pt minus 2pt}%
\setlength{\belowdisplayskip}{4pt plus 2pt minus 2pt}%

\VerdierSemanticAnchor{ch3d-1-3-4-c}
\noindent\textit{c) Every distinguished triangle in $D(\mathcal A)$ is
isomorphic to a distinguished triangle of the form
$(X,Y,Z,[u],[v],\delta(S))$, where
$S=0\longrightarrow X\xrightarrow{u}Y\xrightarrow{v}Z\longrightarrow0$
is an exact sequence in $\mathcal A$.}

\VerdierSemanticAnchor{ch3d-1-3-4-proof}
The proof of a) and b) is left to the reader. Let us prove c). By (Chap. II,
2.2.6), every distinguished triangle in $D(\mathcal A)$ is isomorphic to
the image under $Q\colon K(\mathcal A)\longrightarrow D(\mathcal A)$ of a
distinguished triangle in $K(\mathcal A)$. By (Chap. I, 3.3.1) and axiom
TRII, every distinguished triangle in $K(\mathcal A)$ is isomorphic to a
triangle of the form:

\VerdierSemanticAnchor{ch3d-1-3-4-cone-triangle}
\[
Y\xrightarrow{\dot{p}(u)}c(\dot u)
\xrightarrow{\dot{q}(u)}X[1]
\xrightarrow{-\dot u[1]}Y[1].
\]
Consequently, every distinguished triangle in $D(\mathcal A)$ is
isomorphic to a triangle of the form:

\VerdierSemanticAnchor{ch3d-1-3-4-derived-triangle}
\[
(Y,c(\dot u),X[1],[p(u)],[q(u)],-[u][1]).
\]
Now observe that

\VerdierSemanticAnchor{ch3d-1-3-4-cone-exact}
\[
S=0\longrightarrow Y\xrightarrow{p(u)}c(u)
\xrightarrow{q(u)}X[1]\longrightarrow0
\]
is an exact sequence in $\operatorname{comp}(\mathcal A)$. It then follows
from (Chap. II, 1.2.4) that there is an isomorphism from the triangle
\[
(Y,c(\dot u),X[1],[p(u)],[q(u)],-[u][1])
\]
to the triangle
\[
(Y,c(\dot u),X[1],[p(u)],[q(u)],\delta(S)).
\]

\VerdierSemanticAnchor{ch3d-1-3-5}
\medskip\noindent\textbf{Proposition 1.3.5.}\quad
\textit{a) Every functor
$F\colon\operatorname{comp}(\mathcal A)\longrightarrow C$ that sends
quasi-isomorphisms to isomorphisms in $C$ factors uniquely through the
functor $\operatorname{comp}(\mathcal A)\longrightarrow D(\mathcal A)$.}

\noindent\textit{b) Every half-exact functor from
$\operatorname{comp}(\mathcal A)$ to an abelian category $B$ that sends
quasi-isomorphisms to isomorphisms factors uniquely through the canonical
$\delta$-functor
$\operatorname{comp}(\mathcal A)\longrightarrow D(\mathcal A)$ and a
cohomological functor from $D(\mathcal A)$ to $B$.}

\VerdierSemanticAnchor{ch3d-1-3-5-proof}
Assertion b) follows from assertion a) and the description of the
distinguished triangles in $D(\mathcal A)$ given by (1.3.4). Assertion a)
follows from the definition of $D(\mathcal A)$ (Chap. II, 2.2.6) and the
following lemma.

\VerdierSemanticAnchor{ch3d-1-3-6}
\medskip\noindent\textbf{Lemma 1.3.6.}\quad
\textit{Let $\mathcal A$ be an additive category (not necessarily
abelian), and let $F$ be a functor from $\operatorname{comp}(\mathcal A)$
to a category $C$. The following conditions are equivalent:}
\endgroup
